% ECONOMICS_PROBLEM: {"id": "H26-143", "title": "Efficient tax enforcement", "jel_code": "H26", "source_id": "OP-0477"}
% !TeX program = xelatex
\documentclass[11pt,a4paper]{article}
\usepackage[margin=23mm]{geometry}
\usepackage{fontspec}
\setmainfont{Latin Modern Roman}
\usepackage{amsmath,amssymb,mathtools}
\usepackage{xcolor,enumitem,booktabs,longtable,array,xurl,hyperref}
\definecolor{muted}{HTML}{53636C}
\hypersetup{colorlinks=true,urlcolor=blue,linkcolor=blue}
\setlength{\parindent}{0pt}
\setlength{\parskip}{5pt}
\newcommand{\R}{\mathbb R}
\newcommand{\E}{\mathbb E}
\newcommand{\N}{\mathbb N}
\newcommand{\ind}{\mathbf 1}
\newcommand{\Var}{\operatorname{Var}}
\newcommand{\Cov}{\operatorname{Cov}}
\newcommand{\supp}{\operatorname{supp}}
\DeclareMathOperator*{\argmin}{arg\,min}
\DeclareMathOperator*{\argmax}{arg\,max}
\newcommand{\entrymeta}[3]{{\sffamily\small\color{muted}\textbf{\texttt{#1}}\quad|\quad #2\par #3\par}}
\newcommand{\Problemref}[1]{\texttt{#1}}
\newcommand{\JELref}[2]{\texttt{#2} (JEL \texttt{#1})}
\begin{document}
\section*{Efficient tax enforcement}
\textbf{Problem ID:} \texttt{H26-143}\quad \textbf{JEL:} \texttt{H26}\par
% BEGIN ECONOMICS STATEMENT
\label{problem:OP-0477}
\label{audited:ECON-088}
{\sffamily\small\color{muted}Original subject group: Taxation and social insurance\par}
\label{definitions:problem:30007759}
\textbf{Audit classification:} Background; proposed extension. Slemrod’s 2019 AEA metadata and abstract are verified. The abstract motivates enforcement research but does not state this joint long-run optimization. Added dynamic policy definitions, liability measurement and the distinction between gap and welfare.
\subsection*{Definitions and precise research target}
Fix a population \(i=1,\ldots,n\), horizon \(H\), discount factor \(0<\beta\leq1\), and budget \(B\). A policy \(a\in\mathcal A\) is a measurable dynamic rule assigning audit probabilities, third-party reporting, filing assistance and norm messages from legally observable histories and public randomization. Declare admissible intensities and the entire timing; a vector of four names is insufficient. Under \(a\), define integrable true statutory liability \(L_{it}(a)\), tax attributable to that period remitted by a specified collection deadline \(P_{it}(a)\), and administrative/compliance cost \(C_H(a)\). Distinguish penalties, refunds and delayed payments. Let
\[
 G_H(a)=\mathbb E\sum_{t=0}^H\beta^t\sum_{i=1}^n[L_{it}(a)-P_{it}(a)],
 \qquad g_B=\inf_{a\in\mathcal A:\ \mathbb E C_H(a)\leq B}G_H(a).
\]
If remittances are capped at correctly assessed liability, the gap is nonnegative; otherwise overpayments must be recorded explicitly. Report \(\varepsilon\)-optimal policies, and claim a minimizer only under attainment conditions. Smaller gaps can be caused by a smaller liability base rather than better collection. Thus also report \(\mathbb E\sum\beta^tL_{it}(a)\), receipts, real activity and private compliance costs. A welfare objective using those quantities is a different policy criterion, not an identity with gap minimization.

Let \(P_O\) describe observed filings, interventions and verification audits. A model class \(\mathcal M\) specifies truthful measurement in representative audits or its detection error, audit selection, treatment assignment, interference through taxpayers/advisers and adaptation. Characterize the identified set of the whole response surface \((G_H(a),C_H(a))_{a\in\mathcal A}\). Choosing the observed best arm is not valid if its true liabilities or other policies' response surfaces are unidentified. Extend observation past intervention withdrawal if durability is the target. The review is verified as background; this joint budget-constrained dynamic optimization is editorial.
\subsection*{Variants and assumption changes}
\begin{enumerate}
\item \textbf{Gap versus welfare (editorial).} Minimize unpaid liability or maximize household welfare net of real enforcement costs under a fixed revenue requirement. These objectives need not rank policies identically.
\item \textbf{Independent versus network assignment (editorial).} Individual randomization with no interference identifies individual effects; network spillovers require cluster or saturation designs and an exposure model.
\item \textbf{Static versus adaptive audit rule (editorial).} Compare a fixed audit probability with history-dependent allocation, declaring exploration probabilities and valid inference after adaptation.
\end{enumerate}
\subsection*{References and source audit}
Slemrod (2019), AEA publisher abstract and metadata, Journal of Economic Literature 57(4), 904--954: verified. The exact joint-policy optimization is editorial. Unopened draft/full-text locators are not used as verified evidence.
\begin{enumerate}\raggedright
\item\hypertarget{definitions:source:30007759:1}{} Joel Slemrod. 2019. \emph{Tax Compliance and Enforcement}. Journal of Economic Literature 57(4), 904–954; publisher abstract.
\url{https://www.aeaweb.org/articles?id=10.1257/jel.20181437}.
DOI: \url{https://doi.org/10.1257/jel.20181437}.
\end{enumerate}

\subsection*{Integrated refinement: ECON-088}
{\sffamily\small\color{muted}Joint choice of taxes and enforcement capacity\par}
This refinement adopts the additional source conventions
(page~\pageref{audited:conventions}), subject to its local restrictions.
\textbf{Joint tax schedules and endogenous enforcement capacity.}
Extend the welfare variant to $T_\tau(y)$ and enforcement $e$, allowing true
activity $y_i(\tau,e)$ and reports $\widehat y_i(\tau,e)$ to respond jointly.
For declared social weights and public-service value $V$, consider
\[
\sup_{\tau,e}\left\{\sum_i\omega_iU_i(\tau,e)+V(G)\right\},\qquad
G=\sum_i\mathbb E[T_\tau(\widehat y_i)+\mathrm{penalties}_i]-C(e).
\]
Specify audit/detection, behavior, legal policy constraints, a feasible
nonnegative budget, finite means and equilibrium selection. Taxes and penalties
finance spending and are transfers; do not subtract administration again if its
full welfare effect is already included through $G$. State any additional
unpriced compliance costs separately. Joint welfare optimization differs from
the baseline gap or revenue objective and requires common utility normalization,
real-activity and general-equilibrium incidence responses.
\subsection*{Supplied source audit for ECON-088}
Local [S1], [S2], etc. in this audit and the references immediately below
belong to ECON-088, independently of the original entry's reference numbering.
[S1] and [S2] directly support joint tax design and enforcement. This budget/welfare formulation is editorial and requires explicit taxpayer and public-service accounting.
\subsection*{References for integrated refinement ECON-088}
{\footnotesize\raggedright
\par\noindent\textbf{[S1]} Anders Jensen, Anne Brockmeyer, and Lucie Gadenne (2024). \emph{Taxation and Development}. VoxDevLit 12(1), September 2024.\par \textit{Verified locator / source role:} Primary publisher page checked for review identity, authors, version and background. The specific published agenda is corroborated in [S2]; the exact formal target is editorial.\par \url{https://voxdev.org/voxdevlit/taxation}\par\smallskip
\par\noindent\textbf{[S2]} Oliver Hanney / VoxDev (living compilation). \emph{Evidence gaps: Key unanswered questions in development economics}. Accessed 8 October 2026. The page currently covers 20 topics; the ``16-topics URL is a historical slug.\par \textit{Verified locator / source role:} Topic heading: Taxation and Development. Supports the broad research agenda; this is not a verbatim quotation or an attribution of the displayed mathematical target.\par \url{https://voxdev.org/topic/evidence-gaps-key-unanswered-questions-across-16-topics-development-economics}\par\smallskip
}

\subsection*{Common conventions: Source mathematical conventions}

These conventions apply to each entry unless explicitly replaced.

\textbf{Models and identification.} A primitive model \(m\) specifies all probability laws, feasible choices, information, equilibrium and measurement rules used by its target. The admissible class \(\mathcal M\) is fixed independently of outcomes. If \(P_O(m)\) is its observable probability law and \(\tau(m)\) the target, its identified set at observable law \(P\) is
\[
 \mathcal I_\tau(P)=\{\tau(m):m\in\mathcal M,\ P_O(m)=P\}.
\]
Point identification means that this set is a singleton. An empty set signals incompatibility of data and assumptions. Coordinate bounds need not describe the joint set. Bounds are sharp only if every retained target is attainable or arbitrarily approachable by compatible models. Finite bounds require explicit restrictions on unobserved quantities; unrestricted confounding, hidden wealth, extrapolation or arbitrary equilibrium selection can yield uninformative sets.

\textbf{Causal interventions.} A potential outcome \(Y_i(p)\) is the outcome under a complete policy regime \(p\), including timing, financing, information and market spillovers. The target population and units are fixed before treatment. With interference, use the complete assignment vector or a justified exposure mapping; individual no-interference notation alone cannot identify an economy-wide intervention. A difference of mean potential outcomes does not require the cross-world joint distribution. Conditioning on post-treatment migration, survival, enrollment or employment changes the estimand and needs separate justification.

\textbf{Statistical statements.} An estimator, confidence set, forecast or test specifies a sampling law, model class and loss/coverage criterion. Uniform \((1-\alpha)\)-coverage means \(\inf_{m\in\mathcal M}\Pr_m\{\tau(m)\in C_n\}\geq1-\alpha\), or an explicitly asymptotic version. The whole parameter space trivially has coverage, so informativeness requires a stated width, risk or power objective. A forecasting improvement is relative to a named benchmark, information set and evaluation population.

\textbf{Equilibrium and optimization.} An equilibrium comprises feasible plans, optimality, beliefs where needed, budget constraints and all market/resource clearing. The primitives or a stated benchmark define each correspondence \(\mathcal E(p;m)\). Nonemptiness is assumed or proved, never inferred just from compact policy sets. A selected-equilibrium target uses a declared selection rule; a robust target quantifies over all permitted equilibria. Use suprema or infima unless attainment is established. An \(\varepsilon\)-optimal policy is feasible and within \(\varepsilon\) of the finite optimum value. Report an empty feasible set separately.

\textbf{Welfare and accounting.} Normative utility, interpersonal normalization, social weights, numeraire, discounting and the use of public receipts are primitives. Behavioral utility need not coincide with the welfare criterion. Household welfare includes ownership income and rents once; adding profits again to compensating variation generally double-counts them. A monetary equivalent is defined by an explicit transfer and price convention, with monotonicity and range conditions. Average effects, mechanism contributions, transfers and welfare losses are different targets. Infinite sums require convergence and dynamic feasible plans require terminal or no-Ponzi conditions.

\textbf{Source versions.} A working-paper date, revision date and journal publication year are distinguished. A DOI for a journal version is not automatically the DOI of an earlier manuscript. Locators refer to the stated version and printed pages where specified. A metadata-only check does not verify a claim about a full-text conclusion. Every entry's source subsection narrows the provenance claim accordingly.

\subsection*{Common conventions: Additional-source status and mathematical conventions}

\label{audited:conventions}

Of the 100 supplied ECON entries, 65 distinct problems appear here; 32 close
matches refine earlier entries, and three duplicate questions map to existing
entries. Appendix~\ref{app:audited-crosswalk} lists every destination.
The attachment's P/R/S/U distinctions and source audits are retained. Its
precise mathematical formalizations are editorial unless the individual audit
says otherwise. Candidate subsidy targets ECON-004--006 retain their U flags;
the resolved approval-core question ECON-007 updates OP-0202 above.
The following supplied conventions govern both this part and the attributed
ECON refinements elsewhere in the catalogue, subject to local restrictions.

\textbf{Parameterized questions.} A problem family lists inputs that must be fixed before an instance is evaluated: population, horizon, policies, information, data law, model class and welfare weights. These are required choices, not quantities the citations are assumed to specify. Undefined applications should not be treated as identified estimands.

\textbf{Indices and integrability.} Sets of agents, firms and regions are finite unless a population measure is explicitly used. \(i,j\) index units, \(r\) regions, \(t\) dates and \(h\) horizons. \(\mathbb E\) and \(\Pr\) refer to the declared population and law; \(\mathbf1\{A\}\) is an indicator. Every displayed expectation and discounted sum needs existence in the stated domain. Infinite horizons use a discount in \((0,1)\) and a convergence condition; finite horizons may use discount one.

\textbf{Optimization and existence.} A displayed \(\max\), \(\min\), or \(\arg\max\) is conditional on attainment. For finite feasible menus this is immediate; general spaces need continuity and compactness/coercivity or another existence argument. Without attainment, the target is the supremum/infimum and arbitrarily near-optimal feasible choices. \(\inf\varnothing=+\infty\). Empty feasibility and an unidentified target are different issues.

\textbf{Potential outcomes.} \(Y_i(z)\) is the outcome under a specified intervention, not the observed conditional mean among people choosing \(z\). In binary comparisons \(\Delta Y_i=Y_i(1)-Y_i(0)\). Specify assignment, treatment versions, compliance, information and observation horizon. Interference is allowed under a full policy vector; compressed exposure notation needs a justified mapping. No interference, consistency and overlap are substantive assumptions, not automatic consequences of notation.

\textbf{Populations and observation.} Fix baseline eligibility and population weights. Conditioning on post-policy employment, survival, relocation or program uptake can change the population being compared. Death is not ordinary loss to follow-up; outcomes truncated by death need a composite convention, joint survival target, or explicitly defined survivor stratum. Fertility-changing interventions need a population functional rather than an unexplained fixed descendant cohort. Measurement and missingness models must be supplied.

\textbf{Identification.} A structural model \(m\in\mathcal M\) implies observable law \(P_m\) and target \(\theta(m)\). Identification means
\[
 P_{m_1}=P_{m_2}\ \Longrightarrow\ \theta(m_1)=\theta(m_2)
 \quad\text{for all }m_1,m_2\in\mathcal M .
\]
Without identification, the target is
\[
 \Theta(P;\mathcal M)=\{\theta(m):m\in\mathcal M,\ P_m=P\}.
\]
Writing \(\theta(P)\) before establishing identification can hide incompatible latent models. Confidence sets additionally need a sampling law, confidence level, asymptotic sequence and uniformity class. Point identification, coverage and informative precision are separate properties.

\textbf{Mediation.} Total effects of randomized assignment are distinct from cross-world natural indirect effects. Belief/effort-channel effects require a meaningful mediator intervention and additional measurement, exclusion or mediation assumptions. Randomization of the initial policy alone does not supply those assumptions. Report bounds or interventional alternatives when the stronger target is unsupported.

\textbf{Equilibrium.} A counterfactual \(W(z)\), \(p(z)\), or \(\mu_t^z\) requires individual optimization, feasibility, government/resource budgets, market clearing and state-transition laws. State equilibrium existence and selection, or report ranges over compatible equilibria. Initial-state integration includes subsequent endogenous transitions; current-state integration must use the policy-dependent distribution. Reduced-form invariance after policy is not assumed.

\textbf{Welfare accounts.} Scores, utility, health, dollars and ecological quantities require explicit conversion before addition. Fees, taxes, subsidies, wages and extortion payments are transfers between parties; distributional weights can make them welfare relevant, but their face value is not automatically a resource loss. State ownership, geographic boundaries, foreign-income weights and fiscal finance. Do not add profits to household welfare already including dividends, or count land capitalization and the underlying benefit twice. A benefit-cost expression must distinguish gross benefits from net welfare and subtract every real cost exactly once.

\textbf{Derivatives and risk.} Log derivatives require positive arguments; local responses require admissible perturbations and specified equilibrium branches. Differentiation under expectations needs regularity/domination, and boundaries may allow only one-sided derivatives. Risk uses a specified law; ambiguity uses a declared nonempty law set on a common history space. Adaptive policies must be nonanticipative. A robust criterion is an editorial choice unless attributed explicitly.

\textbf{Scope overlaps.} Allocation, collective choice, contracts, household bargaining and coordinated policy overlap economics with optimization and game theory. They are retained because they appeared in the original catalogue; no claim of a strictly game-theory-free selection is made.
% END ECONOMICS STATEMENT
\end{document}
